\documentclass[10pt,a4paper]{amsart}
\usepackage[margin=19mm]{geometry}
\usepackage{amsmath,amssymb,mathtools}
\usepackage[scr=rsfs,cal=euler]{mathalfa}
\usepackage{xfrac}
\usepackage[unicode,hidelinks,bookmarks=false]{hyperref}
\newcommand{\ie}{i.e.\ }
\newcommand{\cf}{cf.\ }
\newcommand{\resp}{respectively\ }
\newcommand{\Subsection}{\subsection}
\providecommand{\nobreakdash}{\nobreak\hbox{-}}
\setlength{\emergencystretch}{2em}
\multlinegap0pt
\usepackage{bm}
\usepackage{bbm}
\usepackage{dsfont}
\usepackage{xspace}
\usepackage{cancel}
\usepackage{mathtools}
\usepackage{amsthm}
\makeatletter
\@ifundefined{assumption}{\newtheorem{assumption}{Assumption}}{}
\@ifundefined{theorem}{\newtheorem{theorem}{Theorem}}{}
\makeatother
\title[Testing Conditional Stochastic Dominance at Target Points]{Testing Conditional Stochastic Dominance at Target Points}
\author{Federico A. Bugni, Ivan A. Canay,, Deborah Kim}
\date{Source: arXiv:2503.14747v4 (CC BY 4.0)}
\begin{document}
\maketitle

\noindent\emph{Source and licence.} Testing Conditional Stochastic Dominance at Target Points. Version arXiv:2503.14747v4 (CC BY 4.0), distributed under the Creative Commons Attribution 4.0 International licence, as recorded in the arXiv licence metadata for this exact version and shown on the abstract page https://arxiv.org/abs/2503.14747v4. Full rights notice: licenses/arxiv-2503.14747-CC-BY-4.0.txt.\par\smallskip

\noindent\textit{Editorial scope.} Counted unit: the Theorem whose source label is \texttt{thm:limit-Sn-general} in the article (source0.tex lines 1121--1136, source label \texttt{thm:limit-Sn-general}), delivered together with the article's own complete original proof of it (source0.tex lines 1138--1245, one \texttt{proof} environment of 108 lines). No mathematics is written, repaired or re-derived: statement and proof are the article's own text line for line. Required context retained uncounted: ass:zhat (lines 311--313); ass:Z (lines 315--317); ass:cond-continuous (lines 319--326); eq:Sn\_ell (lines 1461--1463). Every definition, display and label of those lines is retained unchanged. Omitted from this package and not redistributed: 1--310, 314--314, 318--318, 327--1120, 1137--1137, 1246--1460, 1464--2003. Nothing inside a retained excerpt is omitted or altered: every retained body is delivered exactly as the article's lines are written, so no omission receipt of any kind accompanies this package. Citations: the retained excerpts carry no citation command at all, so no bibliography entry of the article is transcribed and no reference is redistributed. Reference adaptation: none was needed and none was made; every reference inside the delivered unit resolves inside it, so no unresolved reference or citation is produced. Printed slips kept literal: no printed word, symbol, hypothesis or number of the counted statement or of its proof was altered. Layout and class: the article's own class is \texttt{article} with packages including amsfonts, amsmath, lipsum, amssymb, amsthm, geometry, multirow, booktabs, changepage, xcolor, xr, natbib, latexsym, graphicx; the wrapper is the lane's pinned \texttt{amsart} editorial wrapper at 10pt on A4, which restates the theorem-like environments the retained text needs and adds the article's own macro definitions that the retained text needs (none), with extra wrapper packages (bm, bbm, dsfont, xspace, cancel, mathtools). The delivered numbering is the unit's own. Assets: no retained span contains a figure environment, a graphics-inclusion command, a PDF-page inclusion, a table or a TikZ picture; no image file is copied into this package, the package declares no asset dependency and no image file is redistributed with it. Licence: version arXiv:2503.14747v4 of this article is distributed under the Creative Commons Attribution 4.0 International licence, as recorded in the arXiv licence metadata for this exact version and shown on the abstract page https://arxiv.org/abs/2503.14747v4. A full rights notice is delivered at licenses/arxiv-2503.14747-CC-BY-4.0.txt. Characters: every retained excerpt is pure ASCII, so the delivered engine, XeLaTeX with its default Latin Modern text font, reports no missing character.
\begin{assumption}\label{ass:zhat}
  For each $z\in\mathcal Z$, there exists an $\mathcal A$-measurable estimator $\hat z_n$ such that $\hat z_n \overset{p}{\to} z$.
  \end{assumption}

\begin{assumption}\label{ass:Z}
For any $\varepsilon >0$ and $z\in \mathcal{Z}$, $P\{Z\in (z-\varepsilon ,z+\varepsilon )\}>0$.
  \end{assumption}

\begin{assumption}\label{ass:cond-continuous}
    For any $z\in\mathcal Z$ and any sequence $\{z_k\}$ in the support of $Z$ such that $z_k\to z$,
    \[
    \sup_{t\in\mathbf R}\big|F_Y(t| z_k)-F_Y(t| z)\big|\to 0
    \quad\text{and}\quad
    \sup_{t\in\mathbf R}\big|F_X(t| z_k)-F_X(t| z)\big|\to 0~.
    \]
  \end{assumption}

\begin{equation}\label{eq:Sn_ell}
  S^{\ell}_n = (S^{\ell}_{n,1}, \dots, S^{\ell}_{n,q}) := (Y^{\ell}_{n_y,[1]}, \dots, Y^{\ell}_{n_y,[q_y]}, X^{\ell}_{n_x,[1]}, \dots, X^{\ell}_{n_x,[q_x]})~.
\end{equation}

\begin{theorem}\label{thm:limit-Sn-general}
  Let Assumptions \ref{ass:zhat}, \ref{ass:Z}, and \ref{ass:cond-continuous} hold with $\mathcal{Z}=\{z_{1},z_{2},\dots,z_{L}\}$, and let $\hat S^{\ell}_n$ be defined as in \eqref{eq:Sn_ell} based on $\hat z_{n,\ell}$ instead of $z_{\ell}$, with distance ties broken by any $\mathcal A$-measurable rule. Then,
  \begin{equation}
      (({\hat S_{n}^{1}})^{\prime},\ldots,({\hat S_{n}^{L}})^{\prime})\overset{d}{\rightarrow}
      (({S^{1}})^{\prime},\ldots,({S^{L}})^{\prime})~,
      \label{eq:limit_randomz_L0}
      \end{equation}
  where, for any $(s_{1}^{\ell},\dots,s_{q_y^\ell}^{\ell},s_{q_y^\ell+1}^{\ell},\dots,
  s_{q^\ell}^{\ell})\in\mathbf{R}^{q^\ell}$ for each $\ell=1,\dots,L$ with $q^\ell=q_y^\ell+q_x^\ell$ fixed, $(({S^{1}})^{\prime},\ldots,({S^{L}})^{\prime})$ has the following distribution:
  \begin{equation}
      P\left\{\bigcap_{\ell=1}^{L}\bigcap_{i=1}^{q^{\ell}}\left\{S_{i}^{\ell}\leq s_{i}^{\ell}
      \right\}\right\}~=~\prod_{\ell=1}^{L}\prod_{i=1}^{q_{y}^{\ell}}F_{Y}(s_{i}^{\ell}|z_{\ell})
      \prod_{j=1}^{q_{x}^{\ell}}F_{X}(s_{j+q_{y}^{\ell}}^{\ell}|z_{\ell})~.
      \label{eq:limit-Sn-general_limit_distribution}
  \end{equation}
  \end{theorem}

\begin{proof}
  For each $\ell=1,\dots,L$, let $\hat g_\ell(Z):=|Z-\hat z_{n,\ell}|$, let $M_y^{\ell}$ denote the subset of the indices $i=1,\dots,n_y$ corresponding to the $q_y^{\ell}$ first $\hat g_\ell$-order statistics $(\hat Z^{\ell}_{n_y,(1)},\dots,\hat Z^{\ell}_{n_y,(q_y^{\ell})})$, and let $M_x^{\ell}$ denote the subset of the indices $j=1,\dots,n_x$ corresponding to the $q_x^{\ell}$ first $\hat g_\ell$-order statistics $(\hat Z^{\ell}_{n_x,(1)},\dots,\hat Z^{\ell}_{n_x,(q_x^{\ell})})$. Since each $\hat z_{n,\ell}$ is $\mathcal A$-measurable by Assumption \ref{ass:zhat} and distance ties are broken by an $\mathcal A$-measurable rule, the sets $M_y^{\ell}$, $M_x^{\ell}$ and the orderings within them are $\mathcal A$-measurable. Let $E_n$ denote the following event:
  \begin{equation*}
  E_{n}~=E_{n_y,y}\cap E_{n_x,x}~,~\text{where}~~
  E_{n_y,y}:=\hspace{-6pt}\bigcap_{1\le\ell<\ell'\le L}\hspace{-6pt}
  \left\{M_y^{\ell}\cap M_y^{\ell'}=\emptyset\right\}
  ~\text{and}~
  E_{n_x,x}:=\hspace{-6pt}\bigcap_{1\le\ell<\ell'\le L}\hspace{-6pt}
  \left\{M_x^{\ell}\cap M_x^{\ell'}=\emptyset\right\}.
  \end{equation*}
  In words, $E_n$ means that the subsets of the data used in each of the $L$ tests are pairwise disjoint within each sample. Note that $E_n\in\mathcal A$. We begin by showing that
  \begin{equation}
  P\{E_{n}\}\rightarrow 1~.\label{eq:limit_randomz_L1}
  \end{equation}
  
  Set $\bar\varepsilon:=\frac{1}{4}\min\left\{\left\vert z_{\ell}-z_{\ell'}\right\vert: \ell\not=\ell',~\ell,\ell'=1,\ldots,L\right\}>0$ and fix $\varepsilon\in(0,\bar\varepsilon)$ arbitrarily. Set 
  \[G_{n}:=\left\{\max\left\vert\hat z_{n,\ell}-z_{\ell}\right\vert\leq \varepsilon/2:\ell=1,\ldots,L\right\}~.\]
  For each $\ell=1,\ldots,L$, let $B_{\ell,y}:=\sum_{i=1}^{n_y}I\{|Z_{i}-z_{\ell}|\leq\varepsilon/2\}$ and $\hat B_{\ell,y}:=\sum_{i=1}^{n_y}I\{|Z_{i}-\hat z_{n,\ell}|\leq\varepsilon\}$, and define $B_{\ell,x}$ and $\hat B_{\ell,x}$ analogously for the second sample. Note that
    \begin{equation}
  G_{n}\bigcap\bigcap_{\ell=1}^{L}\{B_{\ell,y}\geq q_{y}^{\ell}\}
  ~\overset{(1)}{\subseteq}~\bigcap_{\ell=1}^{L}\bigcap_{i=1}^{q_y^\ell}
  \{|\hat Z_{n_{y},(i)}^{\ell}-z_{\ell}|\leq 3\varepsilon/2\}
  ~\overset{(2)}{\subseteq}~E_{n_y,y}~,\label{eq:limit_randomz_L2}
  \end{equation}
  where (1) holds by the fact that, for every $\ell=1,\dots,L$, $G_{n}$ and $\{B_{\ell,y}\geq q_{y}^{\ell}\}$ imply that $|\hat z_{n,\ell}-z_{\ell}|\leq \varepsilon/2$ and $|\hat Z_{n_{y},(i)}^{\ell}-\hat z_{n,\ell}|\leq \varepsilon$ for $i=1,\dots,q_y^{\ell}$, and (2) by the fact that for $\ell\not=\ell'$ we have $|z_\ell-z_{\ell'}|\geq 4\bar\varepsilon>3\varepsilon$, so no sample observation can lie within $3\varepsilon/2$ of both $z_\ell$ and $z_{\ell'}$, and therefore $M_y^{\ell}\cap M_y^{\ell'}=\emptyset$ for all $\ell\not=\ell'$, which implies $E_{n_y,y}$.

The same argument applied to the second sample yields the analogue of \eqref{eq:limit_randomz_L2} with $(B_{\ell,x},q_x^\ell,\hat Z_{n_x,(j)}^{\ell},E_{n_x,x})$ in place of $(B_{\ell,y},q_y^\ell,\hat Z_{n_y,(i)}^{\ell},E_{n_y,y})$. From here, we have that \eqref{eq:limit_randomz_L1} follows if we show that
  \begin{align}
  P\{G_{n}^{c}\}&\rightarrow 0~,\label{eq:limit_randomz_L3}\\
  P\{B_{\ell,y}<q_{y}^{\ell}\}\rightarrow 0~\text{ and }~
  P\{B_{\ell,x}<q_{x}^{\ell}\}&\rightarrow 0~\text{ for all }\ell=1,\ldots,L~,
  \label{eq:limit_randomz_L4}
  \end{align}
  since then $P\{E_n^c\}\leq P\{G_n^c\}+\sum_{\ell=1}^{L}\left(P\{B_{\ell,y}<q_y^\ell\} +P\{B_{\ell,x}<q_x^\ell\}\right)\to 0$. Note that this union bound does not rely on the independence of the two samples: $G_n$ may involve the $Z$ observations from both samples, as is the case when $\hat z_{n,\ell}$ is computed from the pooled data.

  Note that \eqref{eq:limit_randomz_L3} holds by Assumption \ref{ass:zhat} and $L$ being finite. For \eqref{eq:limit_randomz_L4}, we only prove the first convergence, as the second is analogous. Since the count $B_{\ell,y}$ is centered at the non-random point $z_\ell$, $B_{\ell,y}\sim Bi(n_y,P\{|Z_{i}-z_{\ell}|\leq\varepsilon/2\})$, where $P\{|Z_{i}-z_{\ell}|\leq\varepsilon/2\}>0$ by Assumption \ref{ass:Z}. By this and $\max_{\ell=1,\ldots,L}q^{\ell}$ being bounded, $\exists N(\varepsilon)$ s.t.\ for all $n_y\geq N(\varepsilon)$,
  \begin{equation}
  0<\frac{1}{2}P\{|Z_{i}-z_{\ell}|\leq\varepsilon/2\}\leq
  P\{|Z_{i}-z_{\ell}|\leq\varepsilon/2\}-\max_{\ell=1,\ldots,L}\frac{q^{\ell}_y}{n_y}~.
  \label{eq:limit_randomz_L4b}
  \end{equation}
  Then, for all $n_y\geq N(\varepsilon)$, we have
  \begin{align*}
  P\{B_{\ell,y}<q_{y}^{\ell}\}
  &=P\left\{{B_{\ell,y}}/{n_y}-P\{|Z_{i}-z_{\ell}|\leq\varepsilon/2\}
  <{q^{\ell}_y}/{n_y}-P\{|Z_{i}-z_{\ell}|\leq\varepsilon/2\}\right\}\\
  &\overset{(1)}{\leq}P\left\{\left\vert{B_{\ell,y}}/{n_y}
  -P\{|Z_{i}-z_{\ell}|\leq\varepsilon/2\}\right\vert
  >\tfrac{1}{2}P\{|Z_{i}-z_{\ell}|\leq\varepsilon/2\}\right\}
  \overset{(2)}{\rightarrow}0~,
  \end{align*}
  as desired, where (1) holds by \eqref{eq:limit_randomz_L4b} and (2) holds by the LLN as $n_y\to\infty$.

  We are now ready to prove \eqref{eq:limit_randomz_L0}. For any $(s_{1}^{\ell},\dots,s_{q_{y}^{\ell}}^{\ell},s_{q_{y}^{\ell}+1}^{\ell},\dots, s_{q^{\ell}}^{\ell})\in\mathbf{R}^{q^{\ell}}$ for each $\ell=1,\dots,L$ with $q^{\ell}=q_{y}^{\ell}+q_{x}^{\ell}$, we have
  \begin{align}
  &P\left\{\bigcap_{\ell=1}^{L}\bigcap_{i=1}^{q^{\ell}}\left\{\hat S_{n,i}^{\ell}\leq
  s_{i}^{\ell}\right\}\right\}\notag\\
  &{\overset{(1)}{=}E\left[\left.P\left\{\bigcap_{\ell=1}^{L}\left\{
  \bigcap_{i=1}^{q_{y}^{\ell}}\left\{\hat Y_{n,\left[i\right]}^{\ell}\leq s_{i}^{\ell}
  \right\}\bigcap\bigcap_{j=1}^{q_{x}^{\ell}}\left\{\hat X_{n,\left[j\right]}^{\ell}\leq
  s_{j+q_{y}^{\ell}}^{\ell}\right\}\right\}\right\vert\mathcal{A}\right\}\right]}\notag\\
  &\overset{(2)}{=}E\left[\left.P\left\{\bigcap_{\ell=1}^{L}\left\{
  \bigcap_{i=1}^{q_{y}^{\ell}}\left\{\hat Y_{n,\left[i\right]}^{\ell}\leq s_{i}^{\ell}
  \right\}\bigcap\bigcap_{j=1}^{q_{x}^{\ell}}\left\{\hat X_{n,\left[j\right]}^{\ell}\leq
  s_{j+q_{y}^{\ell}}^{\ell}\right\}\right\}\right\vert\mathcal{A}\right\}I\{E_{n}\}\right]
  +o_{n}(1)\notag\\
  &\overset{(3)}{=}E\left[\prod_{\ell=1}^{L}\left.P\left\{\left\{
  \bigcap_{i=1}^{q_{y}^{\ell}}\left\{\hat Y_{n,\left[i\right]}^{\ell}\leq s_{i}^{\ell}
  \right\}\bigcap\bigcap_{j=1}^{q_{x}^{\ell}}\left\{\hat X_{n,\left[j\right]}^{\ell}\leq
  s_{j+q_{y}^{\ell}}^{\ell}\right\}\right\}\right\vert\mathcal{A}\right\}I\{E_{n}\}\right]
  +o_{n}(1)\notag\\
  &\overset{(4)}{=}E\left[\prod_{\ell=1}^{L}\left.P\left\{\left\{
  \bigcap_{i=1}^{q_{y}^{\ell}}\left\{\hat Y_{n,\left[i\right]}^{\ell}\leq s_{i}^{\ell}
  \right\}\bigcap\bigcap_{j=1}^{q_{x}^{\ell}}\left\{\hat X_{n,\left[j\right]}^{\ell}\leq
  s_{j+q_{y}^{\ell}}^{\ell}\right\}\right\}\right\vert\mathcal{A}\right\}\right]
  +o_{n}(1)\notag\\
  &\overset{(5)}{=}E\left[\prod_{\ell=1}^{L}\left\{\prod_{i=1}^{q_{y}^{\ell}}
  F_{Y}(s_{i}^{\ell}|\hat Z_{n_{y},\left(i\right)}^{\ell})\prod_{j=1}^{q_{x}^{\ell}}
  F_{X}(s_{j+q_{y}^{\ell}}^{\ell}|\hat Z_{n_{x},\left(j\right)}^{\ell})\right\}\right]+o_{n}(1)~,
  \label{eq:limit_randomz_L5}
  \end{align}
where (1) holds by the LIE, (2) and (4) by \eqref{eq:limit_randomz_L1} and the conditional probability being bounded by one, (3) by Assumption \ref{ass:zhat} and the fact that, conditional on $E_n \in \mathcal{A}$, the $\ell$-th bracket term is a function of $\{Y_i:i\in M_y^{\ell}\}$ and $\{X_j:j\in M_x^{\ell}\}$ only, which are pairwise disjoint observations across $\ell = 1,\dots,L$, and (5) by Assumption \ref{ass:zhat} and the $\mathcal A$-measurable tie-breaking rule, the indices of the selected observations and their ordering by $\hat g_\ell$ are $\mathcal A$-measurable, so conditional on $\mathcal A$ the vector $(\hat Y_{n,[1]}^{\ell},\dots,\hat Y_{n,[q_y^\ell]}^{\ell}, \hat X_{n,[1]}^{\ell},\dots,\hat X_{n,[q_x^\ell]}^{\ell})$ has independent components with conditional conditional CDFs $F_{Y}(\cdot|\hat Z_{n_{y},(i)}^{\ell})$ and $F_{X}(\cdot|\hat Z_{n_{x},(j)}^{\ell})$.

  Next, we show that
  \begin{align}
  \hat Z_{n_{y},\left(i\right)}^{\ell}\overset{p}{\rightarrow}z_{\ell}~\text{ for all }~
  i&=1,\dots,q_{y}^{\ell}\text{ and }~\ell=1,\ldots,L~,\label{eq:limit_randomz_L6}\\
  \hat Z_{n_{x},\left(j\right)}^{\ell}\overset{p}{\rightarrow}z_{\ell}~\text{ for all }~
  j&=1,\dots,q_{x}^{\ell}\text{ and }~\ell=1,\ldots,L~.\label{eq:limit_randomz_L7}
  \end{align}
  We only show \eqref{eq:limit_randomz_L6}, as \eqref{eq:limit_randomz_L7} can be shown analogously. Fix $\varepsilon\in(0,\bar\varepsilon)$ arbitrarily, with $G_n$ and $B_{\ell,y}$ defined as before. By \eqref{eq:limit_randomz_L2}-\eqref{eq:limit_randomz_L4}, we conclude that
  \begin{equation*}
  P\Big\{\bigcap_{\ell=1}^{L}\bigcap_{i=1}^{q_y^\ell}
  \{|\hat Z_{n_{y},\left(i\right)}^{\ell}-z_{\ell}|\leq 3\varepsilon/2\}\Big\}\rightarrow 1~.
  \end{equation*}
  Since the choice of $\varepsilon\in(0,\bar\varepsilon)$ was arbitrary, \eqref{eq:limit_randomz_L6} follows.

  By \eqref{eq:limit_randomz_L6}, \eqref{eq:limit_randomz_L7}, and Assumption \ref{ass:cond-continuous}, $F_{Y}(s_{i}^{\ell}|\hat Z_{n_{y},(i)}^{\ell})\overset{p}{\rightarrow} F_{Y}(s_{i}^{\ell}|z_{\ell})$ and $F_{X}(s_{j+q_{y}^{\ell}}^{\ell}|\hat Z_{n_{x},(j)}^{\ell})\overset{p}{\rightarrow} F_{X}(s_{j+q_{y}^{\ell}}^{\ell}|z_{\ell})$ for every $i$, $j$, $\ell$, and every $(s_{1}^{\ell},\dots,s_{q^{\ell}}^{\ell})$. By this and the fact that the integrand in \eqref{eq:limit_randomz_L5} is bounded by one and converges in probability to a constant, it converges in $L^1$ and so 
  \begin{equation}
  \lim_{n\rightarrow\infty}P\left\{\bigcap_{\ell=1}^{L}\bigcap_{i=1}^{q^{\ell}}\left\{
  \hat S_{n,i}^{\ell}\leq s_{i}^{\ell}\right\}\right\}~
  =~\prod_{\ell=1}^{L}\left\{
  \prod_{i=1}^{q_{y}^{\ell}}F_{Y}(s_{i}^{\ell}|z_{\ell})\prod_{j=1}^{q_{x}^{\ell}}
  F_{X}(s_{j+q_{y}^{\ell}}^{\ell}|z_{\ell})\right\}~.
  \label{eq:limit_randomz_L9}
  \end{equation}
  By definition of convergence in distribution, \eqref{eq:limit_randomz_L9} implies \eqref{eq:limit_randomz_L0}, as desired.
\end{proof}

\end{document}
